% Generated draft: requires mathematical and source review.
\documentclass[11pt]{article}
\usepackage[margin=25mm]{geometry}
\usepackage{fontspec}
\setmainfont{Latin Modern Roman}
\usepackage{amsmath,amssymb,mathtools,unicode-math}
\setmathfont{Latin Modern Math}
\usepackage{enumitem,needspace,xurl,hyperref}
\setlength{\parindent}{0pt}
\setlength{\parskip}{5pt}
\setlength{\emergencystretch}{3em}
\newcommand{\N}{\mathbb{N}}
\newcommand{\Z}{\mathbb{Z}}
\newcommand{\Q}{\mathbb{Q}}
\newcommand{\R}{\mathbb{R}}
\newcommand{\C}{\mathbb{C}}
\newcommand{\F}{\mathbb{F}}
\newcommand{\A}{\mathbb{A}}
\newcommand{\PP}{\mathbb{P}}
\newcommand{\E}{\mathbb{E}}
\newcommand{\eps}{\varepsilon}
\newcommand{\dd}{\,\mathrm{d}}
\begin{document}
\section*{UnsolvedMath ID: 8400037}
\textbf{Source problem code:} catalogue entry
\subsection{Definitions and mathematical statement}
An input is a binary-encoded sextuple \((a,b,c,d,e,f)\in\Z^6\), promised to satisfy
\[
\Delta=b^2-4ac\ne t^2\quad\text{for every }t\in\Z.
\]
The promise includes negative discriminants; no positivity or nonsingularity condition beyond the stated nonsquare promise is added. A yes-instance satisfies
\[
\exists x,y\in\Z,\qquad ax^2+bxy+cy^2+dx+ey+f=0.
\]
The variables are unrestricted integers, not required to be nonnegative or bounded.

The question is NP-hardness of this promise decision problem. Under the conventional polynomial-time many-one interpretation, for every language \(L\in\mathrm{NP}\) there must be a polynomial-time map \(r\) sending every input string to a sextuple satisfying the nonsquare promise, with
\[
w\in L\iff r(w)\text{ has an integral solution}.
\]
In particular, a reduction must preserve the promise for no-instances as well as yes-instances. Membership in NP, or NP-completeness, is not itself requested.

The supplied research summary labels this solved by Manders--Adleman (1978), citing NP-completeness of binary quadratics and the subfamily \(Ax^2+By=C\). But that subfamily has quadratic coefficients \(a=A,b=c=0\), hence \(\Delta=0\), a square, so it is excluded by this problem's promise. The summary's additional claim that their reduction handles nonsquare discriminants is not accompanied by a reduction or precise theorem. Its integer-variable scope must also be checked rather than importing a nonnegative-variable solvability result. Thus the cited subfamily alone does not certify the promised NP-hardness, and the stored solved assessment remains a source-verification issue.
\subsection{Short English statement}
Is unrestricted integer solvability of a binary quadratic equation NP-hard when its quadratic discriminant is promised nonsquare? The supplied cited linear-in-one-variable subfamily has square discriminant and does not by itself resolve the promised problem.
\subsection{Sources}
\begin{itemize}
\item Local UnsolvedMath registry, permanent ID \(8400037\), exact source code \texttt{catalogue entry}, original label O33a.
\item Leonard Adleman and Kevin McCurley, \emph{Open Problems in Number Theory Complexity II} (1994), O33a, source page 22 / LNCS page 312; \url{https://mccurley.org/papers/open.ps.gz}. Recorded access 2026-07-29; extraction recorded as PDF recovered against coauthor PostScript.
\item Manders and Adleman, ``NP-complete decision problems for binary quadratics'' (1978), as supplied in the research summary. No full publication details, variable-domain statement, or nonsquare-discriminant reduction are included.
\end{itemize}
\end{document}
